% !TEX root = ../lab1_report.tex
% ==========================================================
% 实验五：自主探索实验
% ==========================================================
\section{实验五：自主探索实验}

\begin{quote}
\textbf{说明：}该实验允许学生自由发挥，根据个人兴趣自行制定实验方案，方案获审批后方可进行实验。自主探索实验必须具有完整的实验要素。
\end{quote}

\subsection{实验名称}

（根据个人自选课题填写。）

\subsection{实验目的}

（说明本自主探索实验的目的和预期达到的学习目标。建议结合《航空发动机原理》《航空发动机结构》《燃烧学》《叶轮机械原理》等课程所学知识，选择一个感兴趣的切入点进行深入探索。）

\subsection{实验装置}

（列出实验中使用的设备、软件或虚拟系统。可选用的装置包括但不限于：发动机模拟试车系统、三维电子样机、虚拟装配系统、流动显示视频、国家虚拟仿真实验教学平台等。）

\subsection{实验方案}

\subsubsection{实验思路}

（说明选题背景和实验设计思路。建议先进行文献或教材调研，确定探索方向，再设计具体的实验步骤。）

\subsubsection{实验步骤}

\begin{enumerate}
    \item （步骤1）
    \item （步骤2）
    \item （步骤3）
    \item \ldots
\end{enumerate}

\subsection{数据记录}

（记录实验过程中获得的数据，可包含表格、截图等。）

\subsection{数据处理与分析}

（对实验数据进行分析处理，结合理论知识解释实验现象，得出有意义的结论。）

\subsection{结论}

（总结本实验的收获和认识，反思实验设计的不足之处和可能的改进方向。）
